\documentclass[12pt,a4paper]{article}
\usepackage[utf8]{inputenc} % inutile avec XeLaTeX/LuaLaTeX
\usepackage[T1]{fontenc}
\usepackage{amsmath,amssymb,mathrsfs,tikz,times,pifont}
\usepackage{enumitem}
\usepackage{multicol}
\usepackage{lmodern}
\usepackage{array}
\newcommand\circitem[1]{%
\tikz[baseline=(char.base)]{
\node[circle,draw=gray, fill=red!55,
minimum size=1.2em,inner sep=0] (char) {#1};}}
\newcommand\boxitem[1]{%
\tikz[baseline=(char.base)]{
\node[fill=cyan,
minimum size=1.2em,inner sep=0] (char) {#1};}}
\setlist[enumerate,1]{label=\protect\circitem{\arabic*}}
\setlist[enumerate,2]{label=\protect\boxitem{\alph*}}
\everymath{\displaystyle}
\usepackage[left=1cm,right=1cm,top=1cm,bottom=1.7cm]{geometry}
\usepackage[colorlinks=true, linkcolor=blue, urlcolor=blue, citecolor=blue]{hyperref}
\usepackage{array,multirow}
\usepackage[most]{tcolorbox}
\usepackage{varwidth}
\usepackage{float}
\tcbuselibrary{skins,hooks}
\usetikzlibrary{patterns}

\usepackage{pgfplots}
\pgfplotsset{compat=1.15}
\usepackage{mathrsfs}

\usepackage{amsmath}
\usepackage{amsfonts}
\usepackage{amssymb}
\usepackage{tkz-tab}
\usepackage{mdframed}
\usepackage{tikz}
\usetikzlibrary{arrows, shapes.geometric, fit}

\usepackage{xcolor}
\usepackage{tikz}
\usepackage{tkz-tab}
\usepackage{verbatim} % INDISPENSABLE pour l'environnement comment
\usepackage{yhmath}
\usepackage{array,multirow}
\usepackage[most]{tcolorbox}
\usepackage{varwidth}
\usepackage{graphicx}
\usetikzlibrary{shapes.geometric, positioning}
\usepackage{tabularx}
\usepackage{pifont} % Pour les symboles de coches et croix
\usepackage{eso-pic}         % Pour ajouter des éléments en arrière-plan
\usetikzlibrary{angles,quotes,arrows.meta}
% Définition du rouge pour les titres et soulignements
\definecolor{monRouge}{RGB}{180, 0, 0}

\usepackage{graphicx}
\usetikzlibrary{trees} % Nécessaire pour l'arbre
\usetikzlibrary{shapes.geometric, positioning}
\usepackage{tabularx}
\usepackage{pifont} % Pour les symboles de coches et croix
\usepackage{enumitem}
\usepackage{makecell}
% Définition des symboles du tableau
\newcommand{\coche}{{\color{red}\ding{51}}}
\newcommand{\croix}{{\color{red}\ding{55}}}

% Commande pour ajouter du texte en arrière-plan
\AddToShipoutPicture{
    \AtTextCenter{%
        \makebox[0pt]{\rotatebox{45}{\textcolor[gray]{0.6}{\fontsize{5cm}{5cm}\selectfont PGB}}}
    }
}
\renewcommand{\seriesdefault}{b} % rend tout plus foncé

\begin{document}

\tiny

\underline{\textbf{\textcolor{red}{Exemple}}}

On considère l’ensemble E des entiers de 20 à 40. On choisit l’un de ses nombres au hasard.

\begin{itemize}
\item[•] A est l’événement : « le nombre est multiple de 3 »
\item[•] B est l’événement : « le nombre est multiple de 2 »
\item[•] C est l’événement : « le nombre est multiple de 6 ».
\end{itemize}
Calculer $p(A)$, $p(B)$, $p(C)$, $p(A\cap B)$, $p(A\cup B)$, $p(A\cap C)$ et $p(A\cup C)$.

+++++++++++++++++++++++++++++++++++++++++++++++++++++++++++++++++++++++++++++++++++++++++++++++

\underline{\textbf{\textcolor{red}{Solution}}}\\

1. Analyse de l'univers $\Omega$

L'ensemble $E$ contient les entiers de 20 à 40.
$\text{Card}(\Omega) = 40 - 20 + 1 = 21$

2. Probabilités simples

\begin{itemize}
    \item \textbf{Événement A :} Multiples de 3 dans $\{20, \dots, 40\}$
    
    $A = \{21, 24, 27, 30, 33, 36, 39\} \Rightarrow \text{Card}(A) = 7$ donc $ p(A) = \frac{7}{21} = \frac{1}{3}$

    \item \textbf{Événement B :} 
    Multiples de 2 (nombres pairs)
    
    $B = \{20, 22, 24, 26, 28, 30, 32, 34, 36, 38, 40\} \Rightarrow \text{Card}(B) = 11$ donc $ p(B) = \dfrac{11}{21}$

    \item \textbf{Événement C :} 
    Multiples de 6
    
    $C = \{24, 30, 36\} \Rightarrow \text{Card}(C) = 3$ donc $ p(C) = \frac{3}{21} = \frac{1}{7} $
\end{itemize}

3. Probabilités composées

\begin{itemize}
    \item \textbf{Intersection $p(A \cap B)$ :}
    $A \cap B$ sont les multiples de $3 \times 2 = 6$. On constate que $A \cap B = C$. Donc
    $ p(A \cap B) = \frac{3}{21} = \frac{1}{7} $

    \item \textbf{Union $p(A \cup B)$ :}
    $p(A \cup B) = p(A) + p(B) - p(A \cap B)$. Donc 
    $ p(A \cup B) = \frac{7}{21} + \frac{11}{21} - \frac{3}{21} = \frac{15}{21} = \frac{5}{7} $

    \item \textbf{Intersection $p(A \cap C)$ :}
    Comme $C \subset A$ (tout multiple de 6 est multiple de 3), $A \cap C = C$. Donc
    $ p(A \cap C) = \frac{3}{21} = \frac{1}{7} $

    \item \textbf{Union $p(A \cup C)$ :}
    Comme $C \subset A$, $A \cup C = A$. Donc
    $ p(A \cup C) = \frac{7}{21} = \frac{1}{3} $
\end{itemize}

++++++++++++++++++++++++++++++++++++++++++++++++++++++++++++++++++++++++++++++++++++++++++++\\

\underline{\textbf{\textcolor{red}{Exemple}}}\\

On lance un dé parfait de 6 faces numéroté de 1 à 6. 

Calculer la probabilité des éléments suivant : \\
$A$: "Le numéro apparu est pair" \\
$B$: "Le numéro apparu est supérieur à 3" \\
$C$: "Le numéro apparu est un multiple de 3" \\

+++++++++++++++++++++++++++++++++++++++++++++++++++++++++++++++++++++++++++++++++++++++++++++\\

\underline{\textbf{\textcolor{red}{Solution}}}\\

$\Omega = \{1; 2; 3; 4; 5; 6\} \quad A = \{2; 4; 6\} \quad B = \{4; 5; 6\} \quad C = \{3; 6\}$

\begin{enumerate}
    \item $P(A) = \frac{\text{Card } A}{\text{Card } \Omega} = \frac{3}{6} = \frac{1}{2}$
    \item $P(B) = \frac{\text{Card } B}{\text{Card } \Omega} = \frac{3}{6} = \frac{1}{2}$
    \item $P(C) = \frac{\text{Card } C}{\text{Card } \Omega} = \frac{2}{6} = \frac{1}{3}$
\end{enumerate}

+++++++++++++++++++++++++++++++++++++++++++++++++++++++++++++++++++++++++++++++++++++++++\\

\underline{\textbf{\textcolor{red}{Exemple}}}

Une urne contient 3 boules blanches numérotées de 1 à 3, 4 boules rouges numérotées de 1 à 4, 2 boules noires numérotées de 1 à 2. Les boules sont indiscernables au toucher. On tire successivement sans remise trois boules de l'urne. \\
Calculer le nombre de tirage possible. \\
Calculer la probabilité des éléments suivants : \\
$A$: "le tirage est unicolore" \\
$B$: "le tirage contient 2 blanches exactement" \\
$C$: "le tirage contient au moins 1 boule de numéro impair" \\
$D$: "la première boule tirée porte un numéro pair" \\
$E$: "la première boule est noire et la deuxième porte le numéro 2" \\
$F$: "le tirage contient au plus 2 boules rouges"

+++++++++++++++++++++++++++++++++++++++++++++++++++++++++++++++++++++++++++++++++++++++++++++++++

\underline{\textbf{\textcolor{red}{Solution}}}\\
Urne : $\{b_1; b_2; b_3; R_1; R_2; R_3; R_4; n_1; n_2\}$
\begin{enumerate}
    \item $\text{Card } \Omega = A_9^3 = 504$
    \item 
    $\text{Card } A = A_3^3 + A_4^3 = 30 \Rightarrow P(A) = \frac{\text{Card } A}{\text{Card } \Omega} = \frac{30}{504} = \frac{5}{84}$ \\
    
    $\text{Card } B = A_3^2 \times A_6^1 \times \frac{3!}{2!} = 108$ $\Rightarrow$
    $P(B) = \frac{\text{Card } B}{\text{Card } \Omega} = \frac{108}{504} = \frac{3}{14}$ \\
    
    $\text{Card } C = \text{Card } \Omega - \text{Card } \bar{C}$ or
    $\text{Card } \bar{C} = A_4^3 = 24$ donc
    $\text{Card } C = 504 - 24 = 480$ $\Rightarrow$
    $P(C) = \frac{480}{504} = \frac{20}{21}$ 
    
    $\text{Card } D = A_4^1 \times A_8^2 = 224$ $\Rightarrow$
    $P(D) = \frac{\text{Card } D}{\text{Card } \Omega} = \frac{224}{504} = \frac{4}{9}$ \\
    
    $\text{Card } E = A_1^1 \times A_3^1 \times A_7^1 + A_1^1 \times A_2^1 \times A_7^1 = 35$ $\Rightarrow$
    $P(E) = \frac{\text{Card } E}{\text{Card } \Omega} = \frac{35}{504} = \frac{5}{72}$ \\
    
    $\text{Card } F = \text{Card } \Omega - \text{Card } \bar{F}$ or
    $\text{Card } \bar{F} = A_3^3 = 6$ donc
    $\text{Card } F = 504 - 6 = 498$ $\Rightarrow$
    $P(F) = \frac{498}{504} = \frac{83}{84}$
\end{enumerate}

+++++++++++++++++++++++++++++++++++++++++++++++++++++++++++++++++++++++++++++++++++++++++++++++++++
\newpage
\underline{\textbf{\textcolor{red}{Exemple}}}\\
On lance 2 dé parfait de 6 faces numéroté de 1 à 6. On note par $S$ la somme des numéros apparus sur les faces supérieurs des 2 dés.
\begin{enumerate}
    \item À l'aide d'un tableau à double entrée donner toute les valeurs possible de $S$.
    \item Calculer la probabilité des événements suivants :
\end{enumerate}

++++++++++++++++++++++++++++++++++++++++++++++++++++++++++++++++++++++++++++++++++++++++++++

\underline{\textbf{\textcolor{red}{Solution}}}\\
\begin{itemize}
    \item $A = \text{"S est pair"}$
    \item $B = \text{"S} \geqslant 7\text{"}$
    \item $C = \text{"S est un multiple de 3"}$
    \item $D = \text{"S} \geqslant 7 \text{ et S est impair"}$
\end{itemize}

\begin{tabular}{|c|c|c|c|c|c|c|}
\hline
$d_2 \setminus d_1$ & \textbf{1} & \textbf{2} & \textbf{3} & \textbf{4} & \textbf{5} & \textbf{6} \\ \hline
\textbf{1} & 2 & 3 & 4 & 5 & 6 & 7 \\ \hline
\textbf{2} & 3 & 4 & 5 & 6 & 7 & 8 \\ \hline
\textbf{3} & 4 & 5 & 6 & 7 & 8 & 9 \\ \hline
\textbf{4} & 5 & 6 & 7 & 8 & 9 & 10 \\ \hline
\textbf{5} & 6 & 7 & 8 & 9 & 10 & 11 \\ \hline
\textbf{6} & 7 & 8 & 9 & 10 & 11 & 12 \\ \hline
\end{tabular}\\

$\text{Card } \Omega = 36$\\

$
\begin{aligned}
    P(A) &= \dfrac{18}{36} = \dfrac{3}{6} = \dfrac{1}{2} \qquad & P(C) &= \dfrac{12}{36} = \dfrac{2}{6} = \dfrac{1}{3} \\
    P(B) &= \dfrac{21}{36} = \dfrac{7}{12} \qquad & P(D) &= \dfrac{12}{36} = \dfrac{2}{6} = \dfrac{1}{3}
\end{aligned}$

+++++++++++++++++++++++++++++++++++++++++++++++++++++++++++++++++++++++++++++++++++++++++

\underline{\textbf{\textcolor{red}{Exemple}}}\\

Une urne contient trois boules rouges numérotées de 1 à 3, quatre boules vertes numérotées de 1 à 4 et 2 boules noires numérotées de 1 à 2, qui sont indiscernables au toucher. \\
On tire successivement sans remise 3 boules de l'urne.

\begin{enumerate}
    \item Calculer la probabilité des événements suivants :
    \begin{itemize}
        \item $A$ : << les trois boules portent la même couleur >>
        \item $B$ : << la première boule tirée est rouge >>
        \item $C$ : << les boules portent des numéros pairs >>
    \end{itemize}
    \item Calculer $P_B(A)$ et $P_A(B)$.
\end{enumerate}

++++++++++++++++++++++++++++++++++++++++++++++++++++++++++++++++++++++++++++++++++++

\underline{\textbf{\textcolor{red}{Solution}}}\\

$\text{Card } \Omega = A_9^3 = 504$
\begin{enumerate}
\item $P(A) = \dfrac{A_3^3 + A_4^3}{504} = \dfrac{6 + 24}{504} = \dfrac{30}{504} = \dfrac{5}{84}$ \\

$P(B) = \dfrac{A_3^1 \times A_8^2}{504} = \dfrac{3 \times 56}{504} = \dfrac{168}{504} = \dfrac{1}{3}$ \\

$P(C) = \dfrac{A_4^3}{504} = \dfrac{24}{504} = \dfrac{1}{21}$

\item $P(A \cap B) = \dfrac{A_3^3}{504} = \dfrac{6}{504} = \dfrac{1}{84}$ \\

$P_B(A) = \dfrac{P(A \cap B)}{P(B)} = \dfrac{\dfrac{1}{84}}{\dfrac{1}{3}} = \dfrac{3}{84}$ \\
\begin{equation*}
    \boxed{P_B(A) = \dfrac{3}{84} = \dfrac{1}{28}}
\end{equation*}

$P_A(B) = \dfrac{P(A \cap B)}{P(A)} = \dfrac{\dfrac{1}{84}}{\dfrac{5}{84}} = \dfrac{1}{84} \times \dfrac{84}{5} = \dfrac{1}{5}$ \\
\begin{equation*}
    \boxed{P_A(B) = \dfrac{1}{5}}
\end{equation*}

\end{enumerate}
\newpage
++++++++++++++++++++++++++++++++++++++++++++++++++++++++++++++++++++++++++++++

\underline{\textbf{\textcolor{red}{Exemple}}}\\
On dispose de 2 dés non pipés, numérotés de 1 à 6. On lance les deux dés et on fait la somme des numéros apparus sur la face supérieure des 2 dés.

\begin{itemize}
    \item À l'aide d'un tableau à double entrée, déterminer l'ensemble des résultats possibles.
    \item Calculer la probabilité des événements suivants :
    \begin{itemize}
        \item $A$ : << La somme est un multiple de 3 >>
        \item $B$ : << La somme est inférieure ou égale à 10 >>
        \item $C$ : << La somme est impaire >>
    \end{itemize}
    \item Calculer $P_C(A)$, $P_B(A)$ et $P_A(C)$.
\end{itemize}

+++++++++++++++++++++++++++++++++++++++++++++++++++++++++++++++++++++++++++++++++++++++++++++

\underline{\textbf{\textcolor{red}{Solution}}}\\
\begin{enumerate}

\item Tableau des sommes possibles

\begin{center}
\renewcommand{\arraystretch}{1.5}
\begin{tabular}{|c|c|c|c|c|c|c|}
\hline
\text{D1} \(\diagup\) \text{D2} & \textbf{1} & \textbf{2} & \textbf{3} & \textbf{4} & \textbf{5} & \textbf{6} \\  \hline
\textbf{1} & 2 & 3 & 4 & 5 & 6 & 7 \\ \hline
\textbf{2} & 3 & 4 & 5 & 6 & 7 & 8 \\ \hline
\textbf{3} & 4 & 5 & 6 & 7 & 8 & 9 \\ \hline
\textbf{4} & 5 & 6 & 7 & 8 & 9 & 10 \\ \hline
\textbf{5} & 6 & 7 & 8 & 9 & 10 & 11 \\ \hline
\textbf{6} & 7 & 8 & 9 & 10 & 11 & 12 \\ \hline
\end{tabular}
\end{center}
\item Calcul des probabilités
$\text{Card } \Omega = 36$ \\

$P(A) = \dfrac{12}{36} = \dfrac{1}{3}$ \quad (Sommes : 3, 6, 9, 12) \\

$P(B) = \dfrac{33}{36} = \dfrac{11}{12}$ \quad (Toutes sauf 11, 11, 12) \\

$P(C) = \dfrac{18}{36} = \dfrac{1}{2}$ \quad (Sommes impaires)
\item Probabilités conditionnelles
$P_C(A) = \dfrac{\text{card}(A \cap C)}{\text{card } C} = \dfrac{6}{18} = \dfrac{1}{3}$ \\

$P_B(A) = \dfrac{\text{card}(A \cap B)}{\text{card } B} = \dfrac{11}{33} = \dfrac{1}{3}$ \\

$P_A(C) = \dfrac{\text{card}(A \cap C)}{\text{card } A} = \dfrac{6}{12} = \dfrac{1}{2}$

\end{enumerate}
++++++++++++++++++++++++++++++++++++++++++++++++++++++++++++++++++++++++++++++++++

-------------------------------------------------------------------------------------------------------
\newpage
\begin{multicols}{2}
\setlength{\columnseprule}{0.1mm} % La largeur de la ligne verticale entre les colonnes
\underline{\textbf{\textcolor{red}{Exercice}}}\\

Le mois d'Août est très pluvieux. Le premier Août 2020, il a plu. Une étude a montré que :
\begin{itemize}
    \item S'il a plu un jour, la probabilité qu'il pleuve le jour suivant est $1/3$.
    \item S'il n'a pas plu un jour, la probabilité qu'il ne pleuve pas le jour suivant est de $2/5$.
\end{itemize}

On note par :
\begin{itemize}
    \item $A_n$ : "Il a plu au $n$-ième jour du mois d'Août"
    \item et $P_n = P(A_n)$
\end{itemize}

\begin{enumerate}
    \item Donner $P_1$
    \item Calculer $P_2$
    \item Montrer que $P_{n+1} = -\frac{4}{15} P_n + \frac{3}{5}$
    \item On note $u_n = P_n - \frac{9}{19}$
    \begin{enumerate}
        \item Montrer que $(u_n)$ est une suite géométrique.
        \item Exprimer $P_n$ en fonction de $n$.
    \end{enumerate}
    \item Quelle est la probabilité qu'il pleuve le 5 Août ?
\end{enumerate}
\end{multicols}
++++++++++++++++++++++++++++++++++++++++++++++++++++++++++++++++++++++++++++++++++++++

\underline{\textbf{\textcolor{red}{Correction de l'exercice}}}\\

\begin{multicols}{2}
\setlength{\columnseprule}{0.1mm} % La largeur de la ligne verticale entre les colonnes
\underline{\textbf{\textcolor{red}{1. Donner $P_1$}}}\\
L'énoncé précise que le premier Août 2020, il a plu. L'événement $A_1$ est donc certain.
$ P_1 = 1 $

\underline{\textbf{\textcolor{red}{2. Calculer $P_2$}}}\\
D'après la formule des probabilités totales en utilisant le système complet d'événements $\{A_1, \overline{A_1}\}$ :
$ P(A_2) = P(A_2 | A_1) \times P(A_1) + P(A_2 | \overline{A_1}) \times P(\overline{A_1}) $
D'après l'énoncé :
\begin{itemize}
    \item $P(A_2 | A_1) = 1/3$ (probabilité qu'il pleuve s'il a plu la veille)
    \item $P(A_1) = 1$ et donc $P(\overline{A_1}) = 0$
\end{itemize}
$ P_2 = \frac{1}{3} \times 1 + P(A_2 | \overline{A_1}) \times 0 = \frac{1}{3} $

\underline{\textbf{\textcolor{red}{3. Montrer que $P_{n+1} = -\frac{4}{15} P_n + \frac{3}{5}$}}}\\
Appliquons la formule des probabilités totales pour le jour $n+1$ :

$ P_{n+1} = P(A_{n+1} | A_n) P(A_n) + P(A_{n+1} | \overline{A_n}) P(\overline{A_n}) $
On sait que :
\begin{itemize}
    \item $P(A_{n+1} | A_n) = 1/3$
    \item $P(\overline{A_{n+1}} | \overline{A_n}) = 2/5$, donc $P(A_{n+1} | \overline{A_n}) = 1 - 2/5 = 3/5$
    \item $P(\overline{A_n}) = 1 - P_n$
\end{itemize}
En remplaçant :
$ P_{n+1} = \frac{1}{3} P_n + \frac{3}{5} (1 - P_n) \implies$
$ P_{n+1} = \frac{1}{3} P_n + \frac{3}{5} - \frac{3}{5} P_n \implies$
$ P_{n+1} = \left( \frac{5}{15} - \frac{9}{15} \right) P_n + \frac{3}{5} \implies$
$ P_{n+1} = -\frac{4}{15} P_n + \frac{3}{5} $

\underline{\textbf{\textcolor{red}{4. Étude de la suite $(u_n)$}}}\\
Soit $u_n = P_n - \frac{9}{19}$.

\textbf{a) Montrer que $(u_n)$ est une suite géométrique :}
$ u_{n+1} = P_{n+1} - \frac{9}{19} = \left( -\frac{4}{15} P_n + \frac{3}{5} \right) - \frac{9}{19} $
En remplaçant $P_n$ par $u_n + \frac{9}{19}$ :
$ u_{n+1} = -\frac{4}{15} \left( u_n + \frac{9}{19} \right) + \frac{3}{5} - \frac{9}{19} $
$ u_{n+1} = -\frac{4}{15} u_n - \frac{36}{285} + \frac{171}{285} - \frac{135}{285} $
$ u_{n+1} = -\frac{4}{15} u_n $
$(u_n)$ est une suite géométrique de raison $q = -\frac{4}{15}$ et de premier terme $u_1 = P_1 - \frac{9}{19} = 1 - \frac{9}{19} = \frac{10}{19}$.

\textbf{b) Exprimer $P_n$ en fonction de $n$ :}
On a $u_n = u_1 \times q^{n-1}$, soit :
$ u_n = \frac{10}{19} \left( -\frac{4}{15} \right)^{n-1} $
D'où :
$ P_n = \frac{10}{19} \left( -\frac{4}{15} \right)^{n-1} + \frac{9}{19} $

\underline{\textbf{\textcolor{red}{5. Probabilité qu'il pleuve le 5 Août}}}
On calcule $P_5$ en posant $n=5$ :
$ P_5 = \frac{10}{19} \left( -\frac{4}{15} \right)^4 + \frac{9}{19} $
$ P_5 = \frac{10}{19} \times \frac{256}{50625} + \frac{9}{19} \approx 0,4763 $
La probabilité qu'il pleuve le 5 août est d'environ \textbf{0,476} (soit 47,6\%).
\end{multicols}

++++++++++++++++++++++++++++++++++++++++++++++++++++++++++++++++++++++++++++++++++++++++++++++++++++++++++++++
\begin{multicols}{2}
\setlength{\columnseprule}{0.1mm} % La largeur de la ligne verticale entre les colonnes
\underline{\textbf{\textcolor{red}{Exercice}}}\\

Demain sera la rentrée des classes. Fatou est une élève assidue, elle s'est absentée le jour de la rentrée une seule fois sur les 5 dernières années. 
\begin{itemize}
    \item Si elle s'absente un jour, la probabilité qu'elle s'absente le jour suivant est de $1/2$.
    \item Si elle ne s'absente pas un jour, la probabilité qu'elle s'absente le jour suivant est $1/4$.
\end{itemize}

On note $A_n$ : "Fatou s'absente au $n$-ième jour" et $P_n = P(A_n)$.

\begin{enumerate}
    \item Calculer $P_1$
    \item Calculer $P_2$
    \item Montrer que $P_{n+1} = \frac{1}{4} P_n + \frac{1}{4}$
    \item Soit $u_n = P_n - \frac{1}{3}$
    \begin{enumerate}
        \item Montrer que $(u_n)$ est une suite géométrique.
        \item Exprimer $P_n$ en fonction de $n$.
    \end{enumerate}
    \item Au $4^{\text{ème}}$ jour Fatou s'est absentée, quelle est la probabilité qu'elle se fût absentée au $3^{\text{ème}}$ jour ?
\end{enumerate}
\end{multicols}
+++++++++++++++++++++++++++++++++++++++++++++++++++++++++++++++++++++++++++++++++++++++++++++++++++++++++
\begin{multicols}{2}
\setlength{\columnseprule}{0.1mm} % La largeur de la ligne verticale entre les colonnes
\underline{\textbf{\textcolor{red}{Correction de l'exercice}}}\\

\underline{\textbf{\textcolor{red}{1. Calculer $P_1$}}}

D'après les données statistiques de l'énoncé (une absence sur les 5 dernières années le jour de la rentrée) :
$ P_1 = \frac{1}{5} = 0,2 $

\underline{\textbf{\textcolor{red}{2. Calculer $P_2$}}}
En utilisant le système complet d'événements $\{A_1, \overline{A_1}\}$, on a :
$ P(A_2) = P(A_2 | A_1) P(A_1) + P(A_2 | \overline{A_1}) P(\overline{A_1}) $
$ P_2 = \left( \frac{1}{2} \times \frac{1}{5} \right) + \left( \frac{1}{4} \times \frac{4}{5} \right) $
$ P_2 = \frac{1}{10} + \frac{1}{5} = \frac{1}{10} + \frac{2}{10} = \frac{3}{10} $
$ P_2 = 0,3 $

\underline{\textbf{\textcolor{red}{3. Montrer que $P_{n+1} = \frac{1}{4} P_n + \frac{1}{4}$}}}
D'après la formule des probabilités totales pour le passage du jour $n$ au jour $n+1$ :
$ P_{n+1} = P(A_{n+1} | A_n) P(A_n) + P(A_{n+1} | \overline{A_n}) P(\overline{A_n}) $
En remplaçant par les probabilités conditionnelles données :
$ P_{n+1} = \frac{1}{2} P_n + \frac{1}{4} (1 - P_n) $
$ P_{n+1} = \frac{1}{2} P_n + \frac{1}{4} - \frac{1}{4} P_n $
$ P_{n+1} = \frac{1}{4} P_n + \frac{1}{4} $

\underline{\textbf{\textcolor{red}{4. Étude de la suite $(u_n)$}}}
Soit $u_n = P_n - \frac{1}{3}$.

\textbf{a) Montrer que $(u_n)$ est une suite géométrique :}
$ u_{n+1} = P_{n+1} - \frac{1}{3} = \left( \frac{1}{4} P_n + \frac{1}{4} \right) - \frac{1}{3} $
$ u_{n+1} = \frac{1}{4} P_n - \frac{1}{12} = \frac{1}{4} \left( P_n - \frac{1}{3} \right) $
$ u_{n+1} = \frac{1}{4} u_n $
$(u_n)$ est une suite géométrique de raison $q = \frac{1}{4}$ et de premier terme $u_1 = P_1 - \frac{1}{3} = \frac{1}{5} - \frac{1}{3} = -\frac{2}{15}$.

\textbf{b) Exprimer $P_n$ en fonction de $n$ :}
$ u_n = u_1 \times q^{n-1} = -\frac{2}{15} \left( \frac{1}{4} \right)^{n-1} $
D'où :
$ P_n = -\frac{2}{15} \left( \frac{1}{4} \right)^{n-1} + \frac{1}{3} $

\underline{\textbf{\textcolor{red}{5. Probabilité conditionnelle (Probabilité qu'elle se fût absentée au jour 3)}}}
On cherche $P(A_3 | A_4)$. Par la définition de la probabilité conditionnelle :
$ P(A_3 | A_4) = \frac{P(A_4 \cap A_3)}{P(A_4)} = \frac{P(A_4 | A_3) P(A_3)}{P(A_4)} $
Calculons d'abord $P_3$ et $P_4$ :
\begin{itemize}
    \item $P_3 = -\frac{2}{15} (\frac{1}{4})^2 + \frac{1}{3} = -\frac{2}{240} + \frac{80}{240} = \frac{78}{240} = \frac{13}{40}$
    \item $P_4 = \frac{1}{4} P_3 + \frac{1}{4} = \frac{13}{160} + \frac{40}{160} = \frac{53}{160}$
\end{itemize}
$ P(A_3 | A_4) = \frac{\frac{1}{2} \times \frac{13}{40}}{\frac{53}{160}} = \frac{13}{80} \times \frac{160}{53} = \frac{13 \times 2}{53} = \frac{26}{53} $
La probabilité est donc d'environ \textbf{0,49}.

\end{multicols}

+++++++++++++++++++++++++++++++++++++++++++++++++++++++++++++++++++++++++++++++

\underline{\textbf{\textcolor{red}{Exemple BAC 2004}}}


Un porte-monnaie contient quatre pièces de 500 F CFA et six pièces de 200 F CFA. Un enfant tire au hasard et simultanément 3 pièces de ce porte-monnaie.

\begin{enumerate}
    \item Calculer la probabilité de l'événement A : << tirer trois pièces de 500 F >>.
    
    \item Soit $X$ la variable aléatoire égale au nombre de pièces de 500 F figurant parmi les trois pièces tirées.
    \begin{enumerate}
        \item Déterminer la loi de probabilité de $X$.
        \item Calculer l'espérance mathématique et l'écart-type de $X$.
    \end{enumerate}
    
    \item L'enfant répète cinq fois l'expérience en remettant chaque fois les trois pièces tirées dans le porte-monnaie. \\
    Quelle est la probabilité que l'événement A se réalise trois fois à l'issue des cinq tirages ?
\end{enumerate}

+++++++++++++++++++++++++++++++++++++++++++++++++++++++++++++++++++++++++++++++++++++++++++++

\begin{correction}
On a un total de $10$ pièces : $4$ pièces de $500$ F et $6$ pièces de $200$ F.

\bigskip

\textbf{1. Calcul de la probabilité de l'événement $A$ : ``tirer trois pièces de 500 F''}

Le nombre total de tirages possibles de $3$ pièces parmi $10$ est :
\[
\binom{10}{3} = 120
\]

Le nombre de tirages favorables (choisir $3$ pièces parmi les $4$ pièces de $500$ F) est :
\[
\binom{4}{3} = 4
\]

Donc :
\[
P(A) = \frac{\binom{4}{3}}{\binom{10}{3}} = \frac{4}{120} = \frac{1}{30}
\]

\bigskip

\textbf{2. Étude de la variable aléatoire $X$}

$X$ désigne le nombre de pièces de $500$ F obtenues lors du tirage de $3$ pièces.

\medskip

\textbf{(a) Loi de probabilité de $X$}

Les valeurs possibles de $X$ sont : $0, 1, 2, 3$.

\[
P(X = k) = \frac{\binom{4}{k} \binom{6}{3-k}}{\binom{10}{3}}
\]

\[
\begin{array}{|c|c|c|c|c|}
\hline
k & 0 & 1 & 2 & 3 \\
\hline
P_i & \frac{1}{10} & \frac{3}{10} & \frac{1}{10} & \frac{2}{10} \\
\hline
\end{array}
\]
\begin{tabular}{|c|c|}
\hline
$k$ & $P(X = k)$ \\
\hline
0 & $\dfrac{\binom{4}{0}\binom{6}{3}}{120} = \dfrac{20}{120} = \dfrac{1}{6}$ \\
\hline
1 & $\dfrac{\binom{4}{1}\binom{6}{2}}{120} = \dfrac{4 \times 15}{120} = \dfrac{60}{120} = \dfrac{1}{2}$ \\
\hline
2 & $\dfrac{\binom{4}{2}\binom{6}{1}}{120} = \dfrac{6 \times 6}{120} = \dfrac{36}{120} = \dfrac{3}{10}$ \\
\hline
3 & $\dfrac{\binom{4}{3}\binom{6}{0}}{120} = \dfrac{4}{120} = \dfrac{1}{30}$ \\
\hline
\end{tabular}
\bigskip

\textbf{(b) Espérance et écart-type}

\medskip

Espérance :
\[
E(X) = \sum k \cdot P(X=k)
\]
\[
E(X) = 0 \cdot \frac{1}{6} + 1 \cdot \frac{1}{2} + 2 \cdot \frac{3}{10} + 3 \cdot \frac{1}{30}
\]
\[
E(X) = \frac{1}{2} + \frac{6}{10} + \frac{3}{30} = \frac{1}{2} + \frac{3}{5} + \frac{1}{10}
\]
\[
E(X) = \frac{5}{10} + \frac{6}{10} + \frac{1}{10} = \frac{12}{10} = \frac{6}{5}
\]

\medskip

Calcul de $E(X^2)$ :
\[
E(X^2) = 0 + 1^2 \cdot \frac{1}{2} + 2^2 \cdot \frac{3}{10} + 3^2 \cdot \frac{1}{30}
\]
\[
E(X^2) = \frac{1}{2} + \frac{12}{10} + \frac{9}{30}
\]
\[
E(X^2) = \frac{1}{2} + \frac{6}{5} + \frac{3}{10}
\]
\[
E(X^2) = \frac{5}{10} + \frac{12}{10} + \frac{3}{10} = \frac{20}{10} = 2
\]

\medskip

Variance :
\[
V(X) = E(X^2) - (E(X))^2 = 2 - \left(\frac{6}{5}\right)^2
\]
\[
V(X) = 2 - \frac{36}{25} = \frac{50}{25} - \frac{36}{25} = \frac{14}{25}
\]

\medskip

Écart-type :
\[
\sigma(X) = \sqrt{\frac{14}{25}} = \frac{\sqrt{14}}{5}
\]

\bigskip

\textbf{3. Répétition de l'expérience}

On répète $5$ fois l'expérience avec remise. Les tirages sont donc indépendants.

La probabilité de succès (événement $A$) est :
\[
p = \frac{1}{30}
\]

On cherche la probabilité d'obtenir exactement $3$ succès en $5$ répétitions :

\[
P = \binom{5}{3} p^3 (1-p)^2
\]

\[
P = \binom{5}{3} \left(\frac{1}{30}\right)^3 \left(\frac{29}{30}\right)^2
\]

\[
P = 10 \times \frac{1}{27000} \times \frac{841}{900}
\]

\[
P = \frac{8410}{24300000} = \frac{841}{2430000}
\]

\end{correction}

\end{document}
